\chapter{概念界定}

\section{教学技能}

关于教学技能的界定，至今尚无统一的表述。教育学定义：“教学技能是指教师在课堂教学中，基于教学理论，运用专业知识和教学经验等，使学生掌握学科基础知识、基本技能并受到思想教育等所采用的一系列教学行为方式\cite{谢尚果2012论师范生教学技能培养机制的创新}。”美国认知心理学家Gagne等学者将其定义为一种程序性知识，认知策略被视为一种技能，是学习者用来控制自己的心智过程的一系列内源性技能\cite{gagne1984learning}。 Joost Lowyck认为教学技能的定义主要依据两个特征：一方面，教学应分解为一些简单和易于训练的行为成分，另一方面离散的教学行为和技能应与学生的结果有明确的关系\cite{lowyck1980teaching}。

美国斯坦福大学的艾伦教授提出来的十四种分类：变化刺激、导入、结束、非言语性暗示、强化学生参与、流畅的提问、探查性提问、高水平组织的提问、发散性提问、确认、举例说明、演讲、有计划的重复和完整的交流技能\cite{董彦玲2012高师院校数学教师课堂教学技能微格训练研究}。

美国学者莫利按照一定的程序把教学技能分为三大类：教学之前的技能、教学之中的技能和教学之后的技能\cite{董彦玲2012高师院校数学教师课堂教学技能微格训练研究}。

日本东京大学井上光洋提出了：教学设计技能、课堂教学技能（实施、评价、管理、决策等）、学校管理技能、普通教学技能、明确课题实质的教学技能\cite{曹新数学教学技能学习教程}。

综上所述，教师的教学技能包括教学设计、说课、课堂教学、课后辅导、教学评价、教学研究技能等六个部分。



\section{数学教学技能的涵义及分类}

数学教学技能是指数学教师在数学教学过程中，运用数学专业知识和哲学、教育学、心理学等有关教学理论及教学经验，进行数学教学的活动方式。以促进学生的数学学习，从而达到教学目的，是数学课堂上采用了与教师特定的意图有关的意图性行为\cite{董彦玲2012高师院校数学教师课堂教学技能微格训练研究}。不同的学者对数学教学技能的划分也不尽相同。

崔冠之把数学教学技能分为了以下九种：讲解概念和规律的技能；分析讲解例题、习题的技能；利用教学手段的技能；引导学生、启发思考、激发兴趣的技能；使用数学语言的技能；板书的设计和使用的技能；课的引入和结束的技能；课堂教学组织技能\cite{崔冠之1995数学教学技能浅析}。

邵利、罗世敏等把课堂教学技能分为三大类：中学数学课堂教学设计技能、中学数学课堂教学实施技能（导入技能、展开技能、结束技能）、数学课堂教学的媒体运用技能\cite{罗世敏中学数学课堂教学技能实训教程}。

曹新等把数学教学技能分解为三大类：教学前教学技能（教学设计、数学说课）；教学中教学技能（语言、导入、组织教学、讲解、板书、提问、变化、强化、结束）；教学后教学技能（听评课、教学研究）\cite{曹新数学教学技能学习教程}。

综上所述，各位学者关于数学教学技能的分类的观点较多，意见并不完全统一，但通过分析、比较可以得出，数学教学技能主要分为教学设计技能、数学说课技能、课堂教学技能（语言、导入、讲解、板书、提问、多媒体使用、结束技能）、教学评价技能、教学研究技能。


\subsection{教学设计技能}

教学设计技能，是指教师依据学生的认知发展水平和数学课程培养目标，制定具体教学目标\cite{曹新数学教学技能学习教程}，选择教学内容，预设教学过程，为学生的学习创设优良环境的行为方式。


\subsection{数学说课技能}

数学说课技能，是指在教育教学理论的指导下，教师在备课的基础上，面对同行、专家或教研人员，通过口头语言结合辅助手段，阐述某一课题的教学设计，并与听者共同探讨改善和优化教学的教学行为方式\cite{曹新数学教学技能学习教程}。


\subsection{课堂教学技能}

课堂教学技能是指在课堂教学实施与运作中应用的技能。主要包括语言、导入、讲解、板书、提问、多媒体使用和结束技能。

下面通过《圆的标准方程》的教案来展示各项具体的课堂教学技能。

\SetTblrTemplate{head,foot}{empty}
\begin{longtblr}[]{
  width = \textwidth,
  vlines,hlines,
  colspec = {X[1,c]X[4,l]X[1.8,c]},
  cell{1}{1} = {c=3}{l},
  cell{2}{2} = {c},
  row{3,4} = {m},
  % cell{3}{1} = {r=3}{m},
  % cell{6}{1} = {r=5}{m},
  % cell{11}{1} = {r=2}{m},
  % cell{13}{1} = {r=2}{m},
}
  教学资源与工具：板书、PPT &                             &       \\
  教学环节           & 教学内容                        & 教学技能  \\
  复习导入           & {
    问题1：圆的定义是什么？\\
    问题2：在平面执教坐标系中，确定圆需要几个要素？
  }& {①导入技能 \\ ②提问技能} \\
  探索新知           & {
    \begin{varwidth}{\hsize}
      问题 3: 若圆心 $A$ 的坐标为 $(a, b)$, 半径为 $r$, 推导圆 的标准方程。\\
      解: 设圆上任意一点 $M$ 的坐标为 $(x, y)$,
      $\because |H A|=r$, 由两点间距离公式得:
      \[
        \sqrt{(x-a)^{3}+(y-b)^{2}}=r
      \]
      两边平方得 $(x-a)^{2}+(x-b)^{2}=r^{2}$, 即为所求方程。 教师在黑板上画图, 并引导学生推导证明过程, 再使 用 PPT 展示标准图形。
    \end{varwidth}
  }& {
    ①语言技能\\
    ②讲解技能\\
    ③板书技能\\
    ④媒体运用技能
  } \\
  课堂练习           &     {
    \begin{varwidth}{\hsize}
      \vspace*{4pt}
      例 1: 根据方程写出圆心和半径
      \begin{enumerate}[(1)]
        \item $(x-4)^{2}+(y+1)^{2}=6$
        \item $(x+2)^{2}+y^{2}=9$
      \end{enumerate}
      例 2: 写出下列圆的标准方程
      \begin{enumerate}[(1)]
        \item 圆心为 $C(-3,4)$, 半径是 $\sqrt{5}$
        \item 圆心为 $C(4,3)$, 半径是 1
      \end{enumerate}
    \end{varwidth}\vspace*{4pt}
  }                        & {①讲解技能\\ ②板书技能} \\
  课堂小结           & 
    \begin{varwidth}{\hsize}
      \begin{enumerate}[1.]
        \item 圆的标准方程
        \item 圆的标准方程的两要素
      \end{enumerate}
    \end{varwidth}
                  & ①结束技能
\end{longtblr}

\subsection{教学评价技能}

教学评价技能，是指根据数学教学目标，对师生在课堂上的活动及其所产生的各种变化进行价值评判与交流，从而推动教学反思的过程。教学评价的内容包括评教学目标、评教学实施、评教师教学技能、评教学效果几个方面。


\subsection{教学研究技能}

教学研究技能，是指掌握基本的教学科研方法，对学科的发展趋势有所了解，能吸取科研成果，并且具有信息检索的能力。数学研究的基本过程可以概括为：确定研究问题、查阅研究文献、确定研究框架、具体开展研究、呈现研究过程与研究结论。